\documentclass{article}
\usepackage{amsmath}
\usepackage{amssymb}
\usepackage{enumitem}

\begin{document}

\title{Prioritized Progress Report and Status}
\author{}
\date{}
\maketitle

\section{Completed or Covered at Least Once}
You have made significant progress on key areas of data structures and algorithms. Here’s a summary of what you've covered:
\begin{itemize}
    \item \textbf{Linked Lists}: Covered single, double, and circular linked lists.
    \item \textbf{Stacks and Queues}: Implemented and studied fundamental operations.
    \item \textbf{Binary Trees}: Covered binary tree structures and their traversals.
    \item \textbf{Binary Search Trees}: Studied and implemented insertions, deletions, and traversals.
    \item \textbf{Matrix Multiplication}: Covered Strassen’s Algorithm.
    \item \textbf{Sorting Algorithms}: Implemented Quick Sort, Merge Sort, and Heap Sort.
    \item \textbf{Binary Search}: Studied and implemented binary search on sorted arrays.
    \item \textbf{Red-Black Trees}: Covered insertion, deletion, and all balancing operations, with full implementations in C++ (not Rust for some fixups).
    \item \textbf{Graphs}: Implemented adjacency lists and adjacency matrices, with some additional graph work.
    \item \textbf{HashMaps}: Implemented open-addressing hash tables in C and Rust, using double hashing and quadratic probing, with performance analytics.
\end{itemize}

\section{Foundational Work Covered}
\begin{itemize}
    \item \textbf{Automata and Regular Expressions}: Explored foundational concepts, including finite automata and regular expressions.
\end{itemize}

\section{Skimmed but Not Fully Explored (May Revisit Later)}
\begin{itemize}
    \item \textbf{Chapter 5: Probabilities}: Skimmed; probability theory is useful but not immediately essential for your current focus.
    \item \textbf{Chapter 8: Sorting in Linear Time}: Covered some basic ideas of Counting Sort, Radix Sort, and Bucket Sort.
    \item \textbf{Chapter 9: Medians and Order Statistics}: Skimmed; revisiting may be useful for deeper optimization techniques.
\end{itemize}

\section{To-Do List (Prioritized for Progress)}

\subsection{Core Data Structures and Algorithms (Highest Priority)}
\begin{itemize}
    \item \textbf{Breadth-First Search (BFS)}: Next logical step based on your graph work. BFS is crucial for graph traversal and will lead into other algorithms like Dijkstra's.
    \item \textbf{Depth-First Search (DFS)}: Follow BFS with DFS to complete core graph traversal techniques.
    \item \textbf{Dijkstra’s Algorithm}: Build on BFS and DFS to implement this shortest-path algorithm.
    \item \textbf{Kruskal’s and Prim’s Algorithms}: After DFS, tackle these minimum spanning tree algorithms. They pair well with your understanding of graphs and Union-Find.
    \item \textbf{AVL Trees}: You have done red-black trees; AVL trees are next in the balancing tree family. A natural progression for tree structures.
    \item \textbf{Union-Find (Disjoint Sets)}: Essential for graph algorithms like Kruskal’s. Manageable with your current skills.
\end{itemize}

\subsection{Dynamic Programming (Medium Priority)}
\begin{itemize}
    \item \textbf{Dynamic Programming (Knapsack, Fibonacci)}: After graph algorithms, shift focus to dynamic programming, which is key for solving optimization problems. Start with basic problems like the knapsack problem and Fibonacci sequence.
\end{itemize}

\subsection{Additional Data Structures (Medium Priority)}
\begin{itemize}
    \item \textbf{B-Trees}: Important for databases and file systems. A good addition after AVL trees.
    \item \textbf{Tries}: Useful for solving string problems efficiently.
\end{itemize}

\subsection{Advanced Data Structures (Optional but Good to Know)}
\begin{itemize}
    \item \textbf{Fibonacci Heaps}: Useful for optimizing graph algorithms, especially Dijkstra’s.
    \item \textbf{van Emde Boas Trees}: A niche structure; useful in specific cases, but can be deferred for now.
\end{itemize}

\subsection{Advanced Algorithms (Medium Priority)}
\begin{itemize}
    \item \textbf{A* Search Algorithm}: Builds on Dijkstra's algorithm, but with heuristics for better performance in pathfinding.
    \item \textbf{String Matching Algorithms (KMP, Rabin-Karp)}: Important for text search problems, but can be tackled after core graph algorithms and dynamic programming.
\end{itemize}

\section{Overall Recommendation}
\begin{enumerate}[label=\arabic*.]
    \item Continue with \textbf{BFS} and \textbf{DFS}, as they are foundational for graph traversal.
    \item Move to \textbf{Dijkstra’s Algorithm} and then tackle \textbf{Kruskal’s} and \textbf{Prim’s}.
    \item Implement \textbf{AVL Trees} and \textbf{Union-Find} (Disjoint Sets) to strengthen your grasp of balancing and graph-based data structures.
    \item Once those are completed, shift focus to \textbf{Dynamic Programming} (Knapsack, Fibonacci) to deepen your problem-solving skills.
    \item As a lower priority, explore \textbf{B-Trees} and \textbf{Tries} to broaden your data structures toolkit.
    \item Finally, explore advanced structures like \textbf{Fibonacci Heaps} and \textbf{van Emde Boas Trees}, but only if needed for specific applications.
\end{enumerate}

\end{document}
